%==================== 2 REVISÃO DE LITERATURA ====================
\section{Revisão de Literatura}

\subsection{Especificidades da previsão fotovoltaica}

A série de geração fotovoltaica resulta da modulação de uma forçante astronômica determinística — a trajetória solar — por fatores atmosféricos estocásticos, do que decorrem propriedades estatísticas particulares \citep{markovics2022}. Destacam-se quatro: (i) forte autocorrelação, que fundamenta o uso de defasagens e médias móveis como preditores; (ii) sazonalidade múltipla, com superposição dos ciclos diário e anual, parcialmente conhecível \textit{a priori} por meio de variáveis astronômicas; (iii) truncamento diurno, com fração substancial de observações nulas no período noturno, o que torna a distribuição do alvo fortemente assimétrica \citep{andrade2023}; e (iv) não-estacionariedade, induzida por degradação dos módulos, \textit{soiling} e falhas de equipamento, o que exige procedimentos de validação que respeitem a ordem cronológica.

\subsection{Métodos e horizontes de previsão}

Adota-se a partição parcimoniosa de \citet{markovics2022}, que agrupa as abordagens em físicas, estatísticas e híbridas, compreendendo os métodos de aprendizado de máquina como subconjunto das estatísticas (orientadas por dados). Os métodos físicos baseiam-se em cadeias de modelos que descrevem analiticamente a conversão da irradiância em potência, dispensando dados históricos, mas com acurácia tipicamente inferior à das abordagens orientadas por dados \citep{andrade2023}. Os métodos estatísticos estimam, a partir de dados históricos, uma função que minimiza uma perda sobre o conjunto de treinamento, sendo as técnicas de aprendizado de máquina as mais empregadas e, frequentemente, as de melhor desempenho \citep{markovics2022}.

Quanto ao horizonte, a literatura distingue quatro faixas: altíssimo curto prazo (\textit{nowcasting}), curto prazo (intra-dia), um dia à frente (\textit{day-ahead}) e médio a longo prazo \citep{tian2023}. A previsão \textit{day-ahead} (24\,--\,48\,h à frente) é a de maior relevância para a programação do despacho e a comercialização \citep{markovics2022}. Nessa faixa, distinguem-se duas estratégias: a baseada em previsão numérica do tempo (NWP) e a puramente endógena, que projeta a geração do dia seguinte exclusivamente a partir do histórico observado na usina, explorando a autocorrelação e os padrões cíclicos \citep{andrade2023}. A abordagem endógena, adotada neste trabalho, dispensa o acesso a previsões NWP ao custo de acurácia tipicamente inferior.

\subsection{Modelos de referência e avaliação}

A literatura recomenda contextualizar o desempenho por comparação com um modelo de referência (\textit{baseline}) \citep{sarmas2022}, formalizada pelo \textit{forecast skill score}, $\mathrm{skill} = 1 - \mathrm{RMSE}_{\text{modelo}}/\mathrm{RMSE}_{\text{ref}}$. A persistência ingênua é inadequada para a geração fotovoltaica em horizontes de horas, por ignorar o ciclo diário; recomenda-se, em seu lugar, a persistência inteligente (\textit{smart persistence}), que incorpora um modelo de céu claro e constitui a referência canônica da área \citep{sarmas2022,markovics2022}. Um modelo estatístico linear é usualmente acrescentado como referência de complexidade intermediária, isolando o ganho atribuível à modelagem não linear.

\subsection{Aprendizado de máquina para séries temporais}

O ajuste de um modelo consiste em selecionar, em uma classe de funções $\mathcal{F}$, aquela que minimiza uma função de perda sobre o treino; o objetivo último, porém, é minimizar o risco esperado sobre dados não vistos. No caso de séries temporais, a hipótese de observações independentes e identicamente distribuídas é violada, e a aproximação só se transfere ao futuro na medida em que a estrutura estatística permaneça estável. A formulação supervisionada é tipicamente obtida por janela deslizante, em que o vetor de preditores é composto por valores defasados e covariáveis disponíveis até o instante de emissão, sob a exigência crítica de que nenhuma informação posterior integre esse vetor, sob pena de vazamento de informação (\textit{data leakage}). O controle do compromisso viés--variância é exercido sobretudo por regularização, que penaliza a complexidade do modelo. Por fim, a estimativa de desempenho deve empregar validação cruzada adaptada — \textit{walk-forward} ou janela deslizante — que preserve a causalidade temporal, evitando o vazamento que o \textit{K-fold} convencional, com particionamento aleatório, introduziria \citep{kappler2024}.